\def\examid{26101200004}
\begin{document}
\begin{enumerate}
\item \multipart{1}{8}
The data set \code{card} contains data on young men in the US from the 1976 National Longitudinal Survey of Young Men. You can load the data into R using the \code{wooldridge} library using:

\medskip
\code{library(wooldridge)}\\
\code{dat = wooldridge::card}
\medskip

The main outcome of interest is the variable \textit{lwage}, which is the log of the hourly wage.\\
What is the mean of the variable \textit{lwage} in the dataset?
\begin{opts}
\item 6.2004
\item 6.3520
\item 6.1395
\item 5.9810
\item \ca{None of the above}
\end{opts}

\item How many missing values or NA's are there for the variable \textit{IQ} in the dataset?
\begin{opts}
\item 47
\item \ca{949}
\item 690
\item 0
\item None of the above
\end{opts}

\item Consider the following linear regression model.
\[\begin{aligned}
lwage = \beta_0 + \beta_1(educ) + \beta_2(exper) + \beta_3(exper^2) + \beta_4(black) &+ \\
\beta_5(smsa) + \beta_6(south) + \beta_7(momdad14) + \beta_8(KWW) &+ \varepsilon
\end{aligned}\]
The model includes controls for education, experience and experience squared, and a test score on knowledge of the world of work (\textit{KWW}). Finally, (\textit{black, smsa, south, momdad14}), indicator variables for race, living in a metropolitan area, living in the south, and living with both parents at age 14 are included in the model.\\
Estimate the model above with OLS. What is the estimated coefficient on education?
\begin{opts}
\item 0.05482
\item 0.06213
\item \ca{0.05570}
\item 0.04911
\item 0.05823
\end{opts}

\item What is the sum of squared errors of the model?
\begin{opts}
\item 392.3281
\item \ca{405.3401}
\item 0.3704
\item 410.1124
\item None of the above
\end{opts}

\item Suppose we wanted to create a confidence interval for one of the parameters of the model. If we wanted to create a 90\% confidence interval, what critical value would we use?
\begin{opts}
\item 1.6449
\item 1.2818
\item 1.9608
\item \ca{1.6454}
\item None of the above
\end{opts}

\item Consider the null hypothesis that $H_0:\beta_{momdad14}=0$. What is the strongest conclusion we can draw about this null hypothesis?
\begin{opts}
\item Reject the null at the 0.1\% level
\item Reject the null at the 1\% level
\item Reject the null at the 5\% level
\item Reject the null at the 10\% level
\item \ca{Fail to reject the null at the 10\% level}
\end{opts}

\item Consider the null hypothesis that $H_0:\beta_{educ}=0.05$. What is the p-value associated with this hypothesis test?
\begin{opts}
\item 0.0925
\item 0.8150
\item 0.2533
\item \ca{0.1850}
\item None of the above
\end{opts}

\item What is the covariance between the estimate of the parameter on \textit{educ} and the parameter on \textit{KWW}?
\begin{opts}
\item $9.3490\times10^{-6}$
\item \ca{$-2.7081\times10^{-6}$}
\item $2.7081\times10^{-6}$
\item $-2.6182\times10^{-6}$
\item $-1.2120\times10^{-5}$
\end{opts}

\item \multipart{9}{15}
The data set \code{mroz} contains data on married women in the US from the 1975 Panel Study of Income Dynamics. You can load the data into R using the \code{wooldridge} library using:

\medskip
\code{library(wooldridge)}\\
\code{dat = wooldridge::mroz}
\medskip

We would like to study what are the main determinants of earning a high wage using a linear probability model. To answer this question you will need to add a new variable to the data set called \textit{highwage} which should equal 1 if the woman's hourly wage is above \$5 (i.e., \textit{wage} $>5$) and zero otherwise. Women who did not work have a missing wage, so \textit{highwage} should also be missing for them.\\
Among women with a reported wage, what percent are \textit{highwage}?
\begin{opts}
\item 13.94\%
\item 75.47\%
\item \ca{24.53\%}
\item 56.84\%
\item None of the above
\end{opts}

\item Estimate a linear probability model on \textit{highwage}.\\
Include in your model controls for years of education (\textit{educ}), years of work experience (\textit{exper}), \textit{age}, and an indicator for living in a city (\textit{city}). Your model should have 5 parameters including the intercept.\\
How many observations are used to estimate the model?
\begin{opts}
\item 753
\item \ca{428}
\item 325
\item 423
\item None of the above
\end{opts}

\item What is the correct interpretation of the parameter on \textit{educ}? Holding fixed the variables in the model, one additional year of education increases the probability of earning a high wage by \dots
\begin{opts}
\item 7.7463 percent
\item 0.0775 percentage points
\item 0.7746 percentage points
\item \ca{7.7463 percentage points}
\item None of the above
\end{opts}

\item What is the estimate of the lower bound of the 90\% confidence interval for the intercept of the model?
\begin{opts}
\item $-1.1335$
\item $-0.5772$
\item \ca{$-1.0847$}
\item $-1.0613$
\item None of the above
\end{opts}

\item Consider the null hypothesis that $H_0:\beta_{city}=0$. What is the strongest conclusion we can draw about this null hypothesis?
\begin{opts}
\item Reject the null at the 0.1\% level
\item Reject the null at the 1\% level
\item Reject the null at the 5\% level
\item Reject the null at the 10\% level
\item \ca{Fail to reject the null at the 10\% level}
\end{opts}

\item Re-estimate the model adding husband's years of education (\textit{huseduc}) as an additional control. What is the percentage change in the parameter on \textit{educ} when we control for husband's education?
\begin{opts}
\item $-8.63\%$
\item \ca{8.63\%}
\item 7.95\%
\item 108.63\%
\item None of the above
\end{opts}

\item Re-estimate the model adding husband's years of education (\textit{huseduc}) as an additional control. What is the variance inflation factor for the parameter on \textit{educ} when we control for husband's education? (note: if you are calculating the variance inflation factor from the formula you need to ensure that the same observations are used in the different regressions)
\begin{opts}
\item 1.6279
\item 1.6137
\item 0.6441
\item \ca{1.5525}
\item None of the above
\end{opts}
\end{enumerate}

\keyheader
\begin{keytable}
\krow{1}{e}{6.2618}{\code{mean(dat\$lwage)}; not listed}
\krow{2}{b}{949}{\code{sum(is.na(dat\$IQ))}; 47 = \textit{KWW}, 690 = \textit{fatheduc}}
\krow{3}{c}{0.05570}{lm drops the 47 rows with missing \textit{KWW}, $n=2963$}
\krow{4}{b}{405.3401}{\code{sum(resid(m)\^{}2)}; 0.3704 is $\hat\sigma$}
\krow{5}{d}{1.6454}{\code{qt(0.95, df.residual(m))}, df $=2954$; 1.6449 is \code{qnorm}, 1.2818 uses 0.90}
\krow{6}{e}{$p=0.1134$}{$p>0.10$}
\krow{7}{d}{0.1850}{$t=(0.05570-0.05)/0.004298=1.326$; 0.0925 is one-sided}
\krow{8}{b}{$-2.71\times10^{-6}$}{\code{vcov(m)["educ","KWW"]}; $9.35\times10^{-6}$ is (\textit{educ}, \textit{exper})}
\krow{9}{c}{24.53\%}{\code{mean(dat\$highwage, na.rm=TRUE)}; 13.94\% divides by all 753 women}
\krow{10}{b}{428}{\code{nobs(m1)}: the 325 non-workers have NA outcome and are dropped}
\krow{11}{d}{0.0775}{LPM: $\times100$ = percentage points}
\krow{12}{c}{$-1.0847$}{\code{confint(m1, level=0.90)}; $-1.1335$ is the 95\% bound}
\krow{13}{e}{$p=0.2229$}{$p>0.10$}
\krow{14}{b}{8.63\%}{$100(0.084150/0.077463-1)$}
\krow{15}{d}{1.5525}{\code{car::vif(m2)}. By hand on all 753 rows (wrong sample) gives 1.6279; 1.6137 is the VIF of \textit{huseduc}}
\end{keytable}

\subsection*{R code}
\lstinputlisting{../key_practiceC.r}
\end{document}
